\documentclass{ou-report-vaf}

% Dit template is gemaakt door P.J. Molijn in het kader van zijn afstuderen aan de OU in 2014.
% Waarvoor hartelijk dank.
% Minieme maar belangrijke wijzigingen zijn aangebracht door E.M. van Doorn
% Het template is versimpeld door Sylvia Stuurman, 2019.
% Het template is aangepast voor de ba informatica door Harrie Passier, Tanja Vos en Pekka Aho 2020 
% Het tamplate is aangepast met nieuwe huisstyle door Rick Neeft, 2021

\def\mytitle{How can automated comparison of inferred models help testers finding bugs?}
\def\myauthor{Rick Neeft}
\def\myStudentId{851829973}
\def\myPresentationDate{T.B.D. 2022}

\begin{document}

    \pagenumbering{roman}

    % Elements of the thesis
    %%%% TITLE PAGE %%%%%%%
% front cover should be a empty as possible. Title (subtitle) - Name - student Id and thesis presentation date

%to prevent that the title page will be referred as page 1, 
%which will give the warning that there is a page 1 twice.

\pagestyle{plain}

\begin{titlepage}
    \newgeometry{right=100pt,left=58.8pt}

    \pagecolor{ou-light-gray}
    \afterpage{\nopagecolor}
    
    %% Insert the OU logo at the top right corner of the page
    %% PowerPoint contains 84 pt margin in picture
    \begin{tikzpicture}[remember picture,overlay]
                \node at (current page.north east)[anchor=north east,inner sep=58.8pt]{
            \includegraphics[scale=0.5]{images/ou-text-logo.png}
        };
    \end{tikzpicture}
    
    %% Extra whitespace at the top.
    \vspace*{5\bigskipamount}
    
    %% add title without hypens
    \nohyphens{{\color{ou-red}\Huge\bf \mytitle}}
     
    \bigskip
    %{\large subtitle, if any}
    
    %\bigskip \bigskip
    %by
    \bigskip \bigskip
    
    {\Large\bf \myauthor}\\
    {\large\myStudentId}
    
    {\myPresentationDate}
    
\end{titlepage}

% reset geometry since we changed them for the title page
\newgeometry{right=0.875in,left=.875in}

%%%% END TITLE PAGE %%%%%%%
    
    % Page 2 must be left blank
    \myemptypage
    
    % Page 3 contains all compulsory data.
    \addcontentsline{toc}{chapter}{Compulsory data}

\myauthor, \myStudentId\\
\textbf{\mytitle}\\
\myPresentationDate
\vspace*{\fill}

Open University of the Netherlands, Faculty of Science\\
Master's Programme in Software Engineering\\
\\
Graduation committee

Chair: Prof. Dr. Tanja E.J. Vos\\
Primary supervisor: Dr. Pekka Aho\\
Seconday supervisor: Fernando Pastor Ricós

Course code: IM9906
\newpage
    
    % Acknowledgements
    %\chapter*{Acknowledgements}
\addcontentsline{toc}{chapter}{Acknowledgements}

I want to thank my supervisors, Pekka Aho, Fernando Pastor Ricós and Tanja Vos, for the opportunity and support of my graduation project. The weekly meeting and feedback were always inspiring and gave me good directions on what to do. I am grateful for allowing me to develop the application in C\#. Besides academic learning, using C\# and Blazor is extremely helpful in my day-to-day career.

Thanks to WithSecure, especially Joona Oikarinen and Tatu Aalto, for their input and feedback regarding the Analysis website and the docker setup.

I want to thank my wife for her support and for allowing me to sit in our home office for countless hours. 

Thank you to my family, and special thanks to my mother, who always supported me and pushed me to achieve my life goals. 

Last but not least, I want to thank my employer, VECOZO, for providing me with time and resources to finish my Masters. 

Rick Neeft\\
\myPresentationDate
    %\myemptypage
    
    % Table of contents
    \begingroup
        \hypersetup{linkcolor=black}
        \tableofcontents
    \endgroup
    \newpage
    
    % Summary
    %\chapter*{Summary}
\addcontentsline{toc}{chapter}{Summary}
lorem ipsum blah blah yadda yadda yadda
    %\newpage
    
    % the content of the thesis
    \pagenumbering{arabic}
    
    
    \chapter{Introduction} \label{introduction}
    
Regression testing is considered a good practice when testing new software versions before being released to the general public. 
However, due to shortened release cycles, the time to market has decreased drastically. As a result, software test teams have less time to test all new software features, let alone try all other elements to prevent unwanted side effects \cite{rapid-release-cycle-issues}.
The research proposed aims to apply change detection between two versions of the GUI of the system under test. 

This section introduces the proposed research discussing the background and context, followed by the research problem, research aims, research questions and why this research is important.

\section{Background}
The TESTAR tool solves a significant obstacle when it comes to testing the GUI. With TESTAR, the tester can automatically start testing the GUI without any upfront scripts. TESTAR automatically generates and executes test sequences based on elements derived from the GUI \cite{VosAho2021}. In recent master graduation assignments, TESTAR has been extended with an inferred model generation module \cite{thesisMulders}, and it became easier to integrate TESTAR in build and release pipelines in a DevOps environment \cite{thesisSlomp}. Those two additions make it easier to run TESTAR upon each source code integration and retrieving an inferred model about the GUI. In a proof of concept, two inferred models are compared to get changes between versions of the SUT \cite{stateDiff}.

Some companies might write down changes of the software in a changelog. However, some changelogs might not be complete and unwanted side effects might be missing entirely. Comparing two inferred models shows all the changes between versions of a SUT, even the unwanted ones.

\section{Research problem}
Although the proof of concept for change detection is a good start, it lacks some basic functionality. For example, it compares calculated abstract hashes with each other. As a result, even the slightest change in GUI results in a removed and added state between versions.

In addition to change detection, the change visualisation does not use the build-in visualisation tool. The proof of concept generates an HTML-based report, making it difficult to see where the change can be found and which steps the tester needs to take to navigate to the changed state in the software.

\section{The aim of the research}
The research aims to apply change detection on the inferred model created with the automated GUI testing tool TESTAR and research whether the outcome is helping testers finding bugs quicker. The main research questions are:

\begin{questions}
    \item How to detect changes between two versions of the SUT?
	\item How to visualise the detected differences to the user? 
\item How to validate the results of \ref{rq:detect-changes} and \ref{rq:diff-visualisation}?
\end{questions}

\section{Scope}
The scope of the research is change detection in GUI inferred models created by TESTAR. Researching how to make models and making significant changes to the creation of models are omitted. However, tweaking the model generation, like adding data that is not saved at the moment, and configuring TESTAR to create a good model, are in scope.

\section{Contribution}
When change detection is available in TESTAR, it becomes easy for testers to run the comparison in a continuous integration environment and receive an overview of the changes found in an updated version of the software.

As a side effect, the comparison and visualisation solutions can run outside the context of TESTAR and can be deployed to a docker environment.

\section{Document outline}
This research proposal is structured as follows. In section \ref{intoduction} the context, aim and objectives for this research proposal are introduced, together with the limitations for the expected outcome. Section \ref{background} describes the background of this proposal and contains knowledge that is available but might not be known by the readers. Section \ref{releatedWork} includes an overview of the material that can be seen as a direct foundation for this proposal or work that influences the expected outcome. Section \ref{questions} formulates the research questions for the thesis. The last section \ref{planning} will outline the chosen approach and planning for the graduation assignment.
    \newpage
    % Reference list
    % keep using the alpha style, this makes it easier for the reader to correlate the used citations
    \bibliographystyle{alpha}
    \bibliography{bibliography}
    
    % List of abbreviations
    
    
    % Appendices

\end{document}